% =====================================================================================
% DATOS DEL ALUMNO Y DEL PROYECTO  -->  ESTE ES EL ÚNICO ARCHIVO DE CONFIGURACIÓN QUE DEBES EDITAR
% (además de escribir tus secciones en la carpeta secciones/)
% Cambia solo lo que está entre llaves { } a la derecha de cada \newcommand.
% El sistema de revisión lee estos datos (título, alumno, tipo de proyecto): no cambies los nombres
% de los comandos (\NombreAlumno, \NombreProyecto, \TipoProyecto, ...).
% =====================================================================================

% ---------- Paquete adicional usado por las tablas de este documento ----------------
% tabularx permite columnas de ancho elastico (tipo X) y evita que las tablas se
% desborden del margen. Se declara aqui porque configuracion/preambulo.tex no se toca.
\usepackage{tabularx}

% ---------- Tipo de proyecto (OBLIGATORIO) ------------------------------------------
% Escribe exactamente  sistema   o   investigacion
%   sistema        -> página web, sistema web, API, sistema de escritorio, app móvil
%   investigacion  -> algoritmo, modelo de aprendizaje de máquina u otro trabajo de investigación
% Según el tipo, la Revisión 2 y la 3 piden secciones distintas (ver GUIA_ALUMNO.md).
\newcommand{\TipoProyecto}{sistema}

% ---------- Nivel de estudios (OBLIGATORIO) ------------------------------------------
% Escribe exactamente  ITIID   o   TSU   (debe coincidir con tu registro).
%   ITIID -> Ingenieria en Tecnologias de la Informacion e Innovacion Digital
%   TSU   -> Tecnico Superior en Desarrollo de Software Multiplataforma
% La portada y las cartas escriben el nombre del programa a partir de este dato:
% NO escribas la carrera a mano en ningun otro lugar.
\newcommand{\NivelEstudios}{ITIID}

% ---------- Estilo de citas y referencias -------------------------------------------
% ieee (por defecto) o apa. Pregunta a tu director cuál debes usar.
\newcommand{\EstiloCitas}{ieee}

% ---------- Alumno -------------------------------------------------------------------
\newcommand{\NombreAlumno}{Joshua André Alvarado Tovar}
\newcommand{\Matricula}{2330242}
% Alumnas: pongan  la  y  a ;  alumnos: pongan  el  y  o
\newcommand{\elolaNombreAlumno}{el}
\newcommand{\OA}{o}

% ---------- Proyecto -----------------------------------------------------------------
\newcommand{\NombreProyecto}{Integración de un sistema para la automatización y trazabilidad del proceso de solicitud, consulta y control estadístico de documentos en el Instituto Registral y Catastral del Estado de Tamaulipas}
% Título para el encabezado de cada página. Usa \\ para partirlo en líneas (máximo 3 líneas).
\newcommand{\NombreProyectoheader}{Integración de un sistema para la automatización y trazabilidad \\del proceso documental en el Instituto Registral y Catastral \\del Estado de Tamaulipas}

% ---------- Organismo receptor (empresa o institución donde realizas la estadía) -----
\newcommand{\organismoreceptor}{Instituto Registral y Catastral del Estado de Tamaulipas}
% Si el organismo es femenino (la universidad, la empresa) pon  la ; si es masculino (el instituto) pon  el
\newcommand{\elolaNombreEmpresa}{el}

% ---------- Asesores y evaluador -----------------------------------------------------
\newcommand{\nasesorinstitucional}{Dr. Said Polanco Martagón}
\newcommand{\nasesorempresaria}{Lic. Edgar Francisco Alvarado Tovar}
\newcommand{\nevalador}{\underline{\hspace{6cm}}}

% ---------- Fechas -------------------------------------------------------------------
\newcommand{\ncuatrimestre}{Agosto-diciembre 2026}
\newcommand{\fechaPortada}{Diciembre de 2026}
\newcommand{\fechacarta}{4 de Diciembre de 2026}
\newcommand{\FechaExposicion}{\underline{\hspace{4cm}}}
\newcommand{\HoraExposionFormatoVenticuatroHoras}{\underline{\hspace{2cm}}}
